\begin{table}[!htbp]
\centering
\caption{Performance on Nonstationary-Demand Instances}
\label{tab:nonstationary_results}

\begin{threeparttable}
\resizebox{\textwidth}{!}{%
\begin{tabular}{l *{2}{ccc} c}
\toprule
\multirow{2}{*}{\textbf{Instance}}
& \multicolumn{3}{c}{\textbf{RA-SAC}}
& \multicolumn{3}{c}{\textbf{Rolling-Horizon DA}}
& \multirow{2}{*}{\textbf{Gap}} \\
\cmidrule(lr){2-4}\cmidrule(lr){5-7}
& Mean & SE & Med.
& Mean & SE & Med.
& \\
\midrule
\texttt{50102914}
& \textbf{35.02} & 0.47 & 31.69
& 68.28 & 1.32 & 58.65
& $-48.7\%$ \\
\texttt{98011894}
& \textbf{1032.95} & 13.71 & 964.54
& 1979.92 & 38.82 & 1550.38
& $-47.8\%$ \\
\texttt{02351RFH-005}
& \textbf{204.40} & 3.56 & 168.65
& 215.05 & 1.74 & 211.54
& $-5.0\%$ \\
\texttt{98011941}
& \textbf{914.19} & 9.81 & 882.08
& 1037.31 & 14.55 & 986.04
& $-11.9\%$ \\
\texttt{02352QGG-004}
& \textbf{1332.85} & 11.50 & 1264.81
& 1468.38 & 14.95 & 1354.50
& $-9.2\%$ \\
\bottomrule
\end{tabular}
}
\begin{tablenotes}
\footnotesize
\item Results report the average controllable cost per week over the 13-week
planning horizon. Both methods are evaluated on the same 1,000 demand paths
for each instance, and the RA-SAC entries use the best formally reassessed
checkpoint. The standard error reflects variation across evaluation paths.
The gap is calculated as
$(\text{RA-SAC}-\text{Rolling-Horizon DA})/\text{Rolling-Horizon DA}$;
therefore, a negative value favors RA-SAC.
\end{tablenotes}
\end{threeparttable}
\end{table}
